%% 4. SCORING SYSTEM
%% ============================================================================

%% 4.1 Overview

The scoring system is the core of reasoning-cores System 2. It analyzes proposed code changes and computes an IMPACT REPORT that determines whether the change should be allowed, blocked, or flagged.

The scoring pipeline consists of:

\begin{enumerate}
    \item Parse both before and after source code using Tree-sitter
    \item Build ASTs and call graphs for both versions
    \item Extract structural metrics from both ASTs
    \item Compute Mamba SSM embeddings for both versions
    \item Calculate delta metrics (8-dimensional risk vector)
    \item Compute coherence delta (embedding distance)
    \item Compute architectural impact score
    \item Apply per-file-kind thresholds
    \item Generate human-readable summary and repair hints
\end{enumerate}

%% 4.2 Risk Vector

The system computes 8 risk dimensions, all based on delta semantics (difference between before and after states):

\begin{table*}[t]
    \centering
    \caption{Risk Vector Dimensions and Normalizers}
    \label{tab:risk_vector}
    \begin{tabular}{llp{6cm}l}
        \toprule
        Dimension & Formula & Description & Normalizer \\\n        \midrule
        cyclomatic & max(0, b_after - b_before) & Change in cyclomatic complexity & 20.0 \\\n        fan_in & max(in-deg(after) - in-deg(before), 0) & Change in in-degree (callers) & 8.0 \\\n        fan_out & max(out-deg(after) - out-deg(before), 0) & Change in out-degree (callees) & 12.0 \\\n        depth & max(d_after - d_before, 0) & Change in nesting depth & 40.0 \\\n        churn & len(symdiff(line_set_before, line_set_after)) & Lines added/removed & 200.0 \\\n        coupling & max(edges_after - edges_before, 0) & Change in total call graph edges & 40.0 \\\n        cohesion & max(lack_cohesion_after - lack_cohesion_before, 0) & Change in cohesion metric & 1.0 \\\n        novelty & 1 - max(cos(emb_before, emb_after), 0) & Semantic divergence from embedding & 1.0 \\\n        \bottomrule
    \end{tabular}
\end{table*}

%% 4.3 Coherence Delta

Coherence delta measures the Euclidean distance between embeddings, normalized by hidden size:

coherence_delta = ||emb(after) - emb(before)||_2 / sqrt(hidden_size)

For cold-start files (empty before_src or <32 characters), coherence_delta is set to 0.0.

%% 4.4 Architectural Impact Score

AIS is computed as the linear remap of cosine similarity between the
L2-normalized before/after embeddings:

  AIS = (cos(emb_before, emb_after) + 1) / 2

so AIS ∈ [0, 1] with 1.0 == identical embeddings and 0.0 == maximally
dissimilar embeddings. See src/s2_core.py: ``ais = (cos + 1.0) / 2.0``.

Earlier revisions of this paper described AIS as a sigmoid-weighted
combination of the risk vector; that description was incorrect and has
been removed. The 8-dim risk vector (cyclomatic, fan_in, fan_out,
depth, churn, coupling, cohesion, novelty) feeds the per-dim ceiling
check separately; AIS is independent of those dimensions.

Note: AIS, coherence_delta (= chord distance = sqrt(2 - 2 cos) on
L2-normalized embeddings), and novelty (= 1 - max(cos, 0)) are
algebraically related — each is a scalar transform of cos. They are
emitted as three fields in ImpactReport for backward compatibility,
not as three independent signals.

%% 4.5 Per-File-Kind Thresholds

Different file types have different acceptable risk profiles. We define per-kind thresholds as shown in Table \ref{tab:thresholds}.

\begin{table*}[t]
    \centering
    \caption{Per-File-Kind Thresholds}
    \label{tab:thresholds}
    \begin{tabular}{lccc}
        \toprule
        File Kind & Coherence Delta (cd) & Architectural Impact Score (ais) & Dim Ceiling \\\n        \midrule
        source_code & 1.5 & 0.4 & 0.9 \\\n        test_code & 2.0 & 0.3 & 0.95 \\\n        plan_md & 3.0 & 0.3 & 1.0 \\\n        doc_md & 3.0 & 0.3 & 1.0 \\\n        config & 1.2 & 0.5 & 0.9 \\\n        \bottomrule
    \end{tabular}
\end{table*}

A block is triggered if:
\begin{itemize}
    \item ais < threshold[kind]
    \item coherence_delta > threshold[kind]
    \item Any risk_dim > dim_ceiling[kind]
\end{itemize}

%% 4.6 Mamba SSM Backbone

We use a Mamba-130M State Space Model as our embedding backbone:

\begin{itemize}
    \item 130M parameters
    \item Hidden size: 768
    \item Memory: ~200MB RAM
    \item Runs entirely on CPU
    \item Pre-trained checkpoint: state-spaces/mamba-130m-hf from HuggingFace
\end{itemize}

The Mamba architecture provides linear scaling with respect to sequence length, making it efficient for code analysis where files can be long.

%% 4.7 Impact Report

The scoring system returns an ImpactReport containing:

\begin{itemize}
    \item architectural_impact_score: float (0-1)
    \item coherence_delta: float
    \item risk_vector: list of 8 floats
    \item regression_detected: bool
    \item human_summary: str
    \item risk_labels: list of 8 strings
    \item cumulative_drift: optional float
    \item cold_start: bool
    \item file_kind: optional str
    \item cd_threshold: optional float
\end{itemize}

%% PLACEHOLDER: Add detailed scoring math derivations
%% PLACEHOLDER: Add calibration methodology
%% PLACEHOLDER: Add comparison with other scoring approaches
